% preamble.tex — 情報系論文風テンプレート用プリアンブル

% MdTex defaults mode では builtin の DEFAULT_LATEX_PREAMBLE が入らないため、
% raw_attribute 等で必要になる互換コマンドを定義する。
\providecommand{\passthrough}[1]{#1}

% 日本語・フォント
% ltjsarticle は LuaTeX-ja 前提のクラス。Pandoc 側も fontspec を読み込むため、
% ここでは重複しやすい luatexja-fontspec を明示ロードしない。
\usepackage[haranoaji]{luatexja-preset}

% 数式・図表・コード
% unicode-math は Pandoc 側で読み込まれることがあるため、amssymb は重複衝突を避けて明示しない。
\usepackage{amsmath}
\usepackage{graphicx}
\usepackage{float}
\usepackage{booktabs}
\usepackage{array}
\usepackage{listings}
\lstset{
  basicstyle=\ttfamily\small,
  numbers=left,
  numberstyle=\scriptsize,
  frame=single,
  breaklines=true,
  columns=fullflexible,
  keepspaces=true
}

% 見出し・キャプション
\usepackage{titlesec}
\titleformat{\section}{\large\bfseries}{\thesection}{0.8em}{}
\titleformat{\subsection}{\normalsize\bfseries}{\thesubsection}{0.8em}{}
\titleformat{\subsubsection}{\normalsize\bfseries}{\thesubsubsection}{0.8em}{}
\titlespacing*{\section}{0pt}{1.2em}{0.5em}
\titlespacing*{\subsection}{0pt}{0.9em}{0.3em}
\titlespacing*{\subsubsection}{0pt}{0.7em}{0.2em}

\usepackage[font=small,labelfont=bf]{caption}
\renewcommand{\figurename}{図}
\renewcommand{\tablename}{表}
\renewcommand{\refname}{参考文献}

% ページ・段落
\pagestyle{plain}
\setlength{\parindent}{1em}
\setlength{\parskip}{0pt}
\linespread{0.98}

% URL / link
\usepackage[unicode,hidelinks]{hyperref}
